%==============================================================================
%  CAPÍTULO V
%==============================================================================
\chapter{Procedimiento concursal de reorganización simplificada de MIPEs}
\label{cap:reorganizacion-simplificada}

\fichaCapitulo{Reestructuración ágil para la pequeña empresa.}{Reemplazo de auditorías,
veedores especializados de Categoría B y el reintento del artículo 286 S.}

%==============================================================================
\bloqueTeorico
%==============================================================================

\subsection{La reducción de barreras procedimentales}
Sobre el diagnóstico legislativo que motivó la creación de esta vía, véase
\textcite{aj-ley21563}; sobre sus limitaciones persistentes,
\textcite{prensa-df-brechas-mypes} y \textcite{prensa-exante-lavin}.
\pendiente{Desarrollar: adaptación de las reglas generales a la realidad de las PYMEs
chilenas.}

\subsection{Sustitución de la auditoría externa por declaración jurada}
\pendiente{Desarrollar la justificación de eficiencia y sus riesgos.}

%==============================================================================
\bloqueNormativo
%==============================================================================

\subsection{Presupuestos de admisibilidad y umbrales de la Ley \Nro 20.416}
\pendiente{Desarrollar.}

\subsection{Nominación exclusiva de veedores de Categoría B}

Efecto sobre honorarios y aranceles fijos.

\subsection{La Protección Financiera Concursal simplificada}

Plazo base de 40 días, prorrogable de forma expedita por 30 días
adicionales.\verificar{Confirmar plazos y requisitos de la prórroga.}

\subsection{El rechazo del acuerdo y el artículo 286 S}

\begin{alerta}[Rectificación de un error difundido]
Es frecuente encontrar en la literatura de divulgación la afirmación de que el
\art{286 S} consagraría un «derecho de reintento» del deudor, esto es, la facultad de
proponer un nuevo acuerdo dentro de quinto día con el apoyo de más del 50\,\% del pasivo.
\destacar{El precepto no dice eso.}
\end{alerta}

El \art{286 S} se titula «Rechazo del Acuerdo» y regula la consecuencia del fracaso de la
propuesta. Su contenido efectivo es el siguiente: si la propuesta es rechazada por no
haberse obtenido el quórum necesario, o porque el deudor no otorgó su consentimiento, y
no estuviere constituida la junta de acreedores, el tribunal \destacar{deberá requerir a
la \superir{} la nominación del liquidador} conforme al \art{37}, dentro de los cinco
días siguientes a esa actuación, acompañando los antecedentes de los tres principales
acreedores que consten en la nómina de créditos reconocidos, excluidas las personas
relacionadas con el deudor. Recibido el certificado de nominación, el tribunal dicta la
resolución de liquidación sin más trámite.

El plazo de cinco días es, por tanto, el que tiene \emph{el tribunal} para requerir la
nominación, y no un plazo concedido al deudor para insistir.

\begin{foco}[Qué significa esto en la estrategia]
El rechazo de la propuesta simplificada conduce a la liquidación de modo prácticamente
automático y sin trámite intermedio. No existe una segunda oportunidad legal. Toda la
negociación debe estar cerrada \emph{antes} de la votación, porque después no hay
mecanismo de recuperación.
\end{foco}

\subsection{Enajenación de activos no esenciales}
\pendiente{Desarrollar el derecho preferente de enajenación de activos no esenciales.}

%==============================================================================
\bloquePractico
%==============================================================================

\subsection{Modelo de propuesta de acuerdo de reorganización simplificada}

\begin{escrito}[Modelo — Propuesta de acuerdo de reorganización simplificada]
\small
\noindent\textsc{S.J.L. en lo Civil de} [\textsc{ciudad}] --- Causa Rol [\textsc{rol}]

\medskip
\noindent [\textsc{nombre del representante legal}], en representación de
[\textsc{razón social de la mipe}], en los autos de reorganización simplificada
caratulados «[\textsc{carátula}]», a US.\ respetuosamente digo:

\medskip
Que, en conformidad a lo dispuesto en la \LIR{}, modificada por la \Lreforma{}, vengo en
presentar formalmente la \textbf{propuesta de acuerdo de reorganización simplificada},
para someterla a la deliberación y votación de los acreedores, sobre la base de las
siguientes estipulaciones:

\begin{enumerate}[leftmargin=1.6em]
  \item \textbf{Remisión y condonación.} Se propone la condonación total de los intereses
  corrientes, moratorios, multas y comisiones devengadas a la fecha de la resolución de
  reorganización.

  \item \textbf{Plazo de gracia.} Se establece un plazo de gracia de doce meses para el
  pago de la primera cuota del capital reestructurado, contado desde la publicación del
  acuerdo firme en el \boletin{}.

  \item \textbf{Plan de amortización.} El capital insoluto verificado se pagará en 48
  cuotas mensuales, iguales y sucesivas, sin recargos, a partir del mes decimotercero.

  \item \textbf{Supervisión del acuerdo.} Se propone la designación de un interventor de
  la nómina de la \superir{}, Categoría B, con atribuciones de inspección mensual de
  flujos.
\end{enumerate}

\medskip
\noindent\textbf{Por tanto}, ruego a US.\ tener por acompañada la propuesta y ordenar su
publicación en el \boletin{}.
\end{escrito}

\begin{alerta}[Adaptación obligatoria del modelo]
Este modelo es un esqueleto. Los porcentajes, plazos y garantías deben ajustarse al flujo
proyectado del deudor y al mapa de mayorías. Una propuesta estándar presentada sin
negociación previa con los acreedores mayoritarios tiene alta probabilidad de rechazo.
\end{alerta}

%==============================================================================
\bloqueJurisprudencial
%==============================================================================

\subsection{Incidente de oposición al acuerdo simplificado}
\pendiente{Sistematizar los primeros fallos posteriores a la entrada en vigencia de la
Ley 21.563.}

\subsection{Efecto del rechazo por falta de quórum}

¿Conduce obligatoriamente a la resolución de liquidación simplificada dictada de oficio,
o cabe privilegiar la prórroga judicial para explorar la venta ordenada de la unidad
económica? \pendiente{Sistematizar criterios.}
